% !PRINTED 441
\weberpage{441}
% running head: § 121. Separation of the opposite classes.

If $k$ is the \gesperrt{principal class}, then
\[
2\omega = \sqrt{D} \text{ or } 2\omega = -1+\sqrt{D},
\]
and since $\sqrt{D}$ is purely imaginary, $k=j(\omega)$ is \gesperrt{real} by § 69 (4). % [sic] The German print writes k=j(\omega), although the introduced class-invariant notation is (k)=j(\omega).
Accordingly there follows from (14) and (15), since we now have $(sk)=(s)$:

\gesperrt{The invariants of opposite classes are conjugate imaginary, the invariants of two-sided classes are real.}

If we return to an arbitrary class $k$ and put $sk$ in place of $k$ in formula (15), it follows that:
\weq{16}{(k) = f_s(sk, -\sqrt{D}),}
or, putting $sk=k'$ and writing $f$ for $f_s$, by (21) % [sic] German print has (21), a forward reference in a passage currently numbered (16)--(17)
\weq{17}{\begin{aligned} (k') &= f(k, \sqrt{D}),\\ (k) &= f(k', -\sqrt{D}), \end{aligned}}
where now $k, k'$ may be any two classes of the determinant $D$.

We are hereby in a position to determine the \gesperrt{Galois} group of the class equation, or at least, to begin with, a group of permutations in which the group of the class equation is contained as a divisor.

Let us assume first that $\sqrt{D}$ has been adjoined to the domain of rationality of the rational numbers, and let $(k), (k'), (k'')\ldots$ denote all the class invariants of the discriminant $D$, while
\[
R(k, k', k'', \ldots)
\]
is any rational function of these quantities. By our theorem this function can be expressed rationally by one of the quantities $(k)$, thus say:
\[
R(k, k', k'', \ldots) = R'(k).
\]
If, further, $(s)$ denotes an arbitrary class invariant of the discriminant $D$, and if the function $R$ has the property of remaining unchanged under all $h$ permutations
\[
\begin{pmatrix} k, & k', & k'', & \ldots \\ sk, & sk', & sk'', & \ldots \end{pmatrix},
\]
whose totality we will denote by $\mathfrak{S}$, then $R'(k)=R'(sk)$, and is therefore \gesperrt{rationally expressible}.

\gesperrt{The Galois group of the class equation after the adjunction of $\sqrt{D}$ is therefore contained in the system $\mathfrak{S}$,}
%CONT>>
